% ============================================================
%  Bioinformatics LAB -- Lab 0 (warm-up)
%  The Chromatin Game: a first look at Hi-C maps
%  Author: Sebastian Korsak, Warsaw University of Technology, MiNI
% ============================================================
\documentclass[11pt]{article}

\usepackage[a4paper,margin=1.8cm]{geometry}
\usepackage[utf8]{inputenc}
\usepackage[T1]{fontenc}
\usepackage{charter}           % clean, readable serif text font
\usepackage{xcolor}
\usepackage{enumitem}
\usepackage{titlesec}
\usepackage{hyperref}
\usepackage{fancyhdr}

% --- one subdued accent colour, echoing the lecture's "beaver" theme ---
\definecolor{accent}{HTML}{7A1F2B}
\definecolor{taskbg}{HTML}{F4EDEE}

\hypersetup{colorlinks=true, urlcolor=accent, linkcolor=accent}

\titleformat{\section}{\large\bfseries\color{accent}}{\thesection}{0.6em}{}
\titlespacing*{\section}{0pt}{6pt}{3pt}
\setlist[enumerate]{leftmargin=1.6em, itemsep=1pt, topsep=1pt}

\pagestyle{fancy}
\fancyhf{}
\renewcommand{\headrulewidth}{0pt}
\renewcommand{\footrulewidth}{0.4pt}
\fancyfoot[L]{\small Bioinformatics LAB -- Warsaw University of Technology (MiNI)}
\fancyfoot[R]{\small Sebastian Korsak}

\newcommand{\key}[1]{\texttt{\textbf{#1}}}

% --- a run-in Task heading that flows straight into its own paragraph ---
\newcounter{tasknum}
\newcommand{\task}[1]{%
	\stepcounter{tasknum}%
	\par\vspace{8pt}\noindent%
	{\color{accent}\textbf{Task \thetasknum.}} \textbf{#1} }

\begin{document}
\thispagestyle{fancy}

\begin{center}
	{\LARGE\bfseries\color{accent} Lab 0 -- The Chromatin Game}\\[2pt]
	{\large A first, hands-on look at Hi-C maps and chromatin folding}\\[2pt]
	{\small Warm-up for Laboratory 1 \quad $\cdot$ \quad Time needed: 20--25 minutes}
\end{center}

\vspace{3pt}
\noindent
Hi-C experiments give an $N\times N$ \textbf{contact map}: entry $(i,j)$ counts how
often loci $i$ and $j$ touch in 3D. This lab has you fold a virtual chromosome
yourself and watch that map emerge, so the idea stops being abstract.

\task{Set up an isolated Python environment and install the game.}%
The game is open-source: \url{https://github.com/SFGLab/TheChromatinGame}.
\begin{center}
\begin{minipage}{0.95\linewidth}
\footnotesize
\begin{verbatim}
git clone https://github.com/SFGLab/TheChromatinGame.git
cd TheChromatinGame
python -m venv chromatin_env        # create the environment
source chromatin_env/bin/activate   # on Windows: chromatin_env\Scripts\activate
pip install -r requirements.txt
python main.py
\end{verbatim}
\end{minipage}
\end{center}
\noindent A dedicated environment per project -- here via \texttt{venv}, or
\texttt{pyenv}/\texttt{conda} if you prefer -- keeps packages from clashing,
a habit worth starting on day one.

\task{Open \textbf{Chromatin MiNI-Lab} from the main menu and freely paint
compartments, tie a loop, and move a slider.}%
Press \key{C} and drag across the colour ribbon to mark beads red (A) or blue
(B); press \key{L} and click two beads (or a cell on the map) to tie a loop
between them; then drag any slider -- e.g.\ temperature or attraction strength
-- under the heatmap. Everything else is explained in-game (press \key{H}).
Watch the 3D structure and the heatmap respond together, live: what you are
doing here is the easy direction, the one physics is good at, where your
choices fix the structure and the structure fixes the map.

\task{Go to \textbf{Single Player}, level 1 (normal mode), and reproduce the
target heatmap shown to you.}%
A target map is displayed next to your own, live one; using the same painting
and loop-tying controls as before, try to make the two look alike, then press
\key{Enter} to freeze the simulation and score your match. Now you are running
the direction biology actually hands you: given the map, recover the structure
that produced it. Many different structures can cast nearly the same shadow,
and noise can point you toward the wrong one -- which is why this is a genuine
\textbf{inverse problem}, not just a puzzle.

\task{Replay the same level in \textbf{hard mode}, where the target's loop
anchors and compartment colours are hidden and only its raw contact map is
shown.}%
Study that map for a moment -- where is the checkerboard strongest, where do
bright squares sit on the diagonal -- before placing any loops or colour, then
score yourself as before. This is the realistic case: a real Hi-C analyst
never sees the ground truth, only the measurement. Normal mode had handed you
the answer key; hard mode does not, exactly like a real dataset on your desk.

\task{Before you leave, write one or two sentences answering: what produced
the checkerboard pattern, what produced the bright square on the diagonal, and
why was hard mode harder?}%
The same logic -- indirect measurement, a physical model, an inferred
structure -- underlies protein structure determination by X-ray, NMR, or
cryo-EM, and it is exactly what we build explicitly, for proteins and
chromatin alike, in the \emph{Molecular Simulations} block later in this
course. This game is your first playable encounter with that idea.

\end{document}
